% GENERATED FILE. Edit canonical lessons, then run npm run generate:books.
% canonical-source-hash: fnv1a64:2a760816eca7e18e
% canonical-lessons: RU-C225-krokodil, RU-C225-zhiraf, RU-C225-zebra, RU-C225-verblyud, RU-C225-olen

\chapter{Crocodile, Giraffe, Zebra, Camel, Deer}
\label{ch:ru-crocodile-giraffe-zebra-camel-deer}
\hlchaptermodality{225}

\begin{chapteropening}
\textbf{By the end of this chapter:} \emph{I can say a crocodile, a giraffe, a zebra, a camel and a deer.}

Complete the last lesson of chapter 225: I can say a crocodile, a giraffe, a zebra, a camel and a deer.
\end{chapteropening}

\section[\texorpdfstring{\ru{крокодил} (krokodíl)}{крокодил (krokodíl)}]{\ru{крокодил} (krokodíl) — a crocodile}

\label{lesson:RU-C225-krokodil}



\begin{quote}
Before the new one: say the Russian for a bush, then the Russian for a sheep.
\end{quote}

\subsection*{You'll want to know: \ru{крокодил}}
\textbf{\ru{крокодил}} (\emph{krokodíl}) — \textquotedblleft{}a crocodile\textquotedblright{}.

\begin{cousinweb}
\textbf{\ru{крокодил}} is a western loan, perhaps through German \emph{Krokodil}, from Latin \emph{crocodīlus}, Greek \emph{krokódeilos}; Old Russian \emph{korkodilŭ} came straight from Greek (Vasmer).
\end{cousinweb}

\begin{grammarlens}[title={What you've built}]
One more word for this chapter.
\end{grammarlens}

\subsection*{Your turn}
\begin{itemize}
\raggedright
  \item \emph{Say it:} \emph{krokodíl}
  \item \emph{Say it:} \emph{krokodíl}, once more
  \item \emph{Say it:} say \emph{krokodíl}
  \item \emph{Recall:} say the Russian for a bush, then the Russian for a sheep, then say \emph{krokodíl} again
\end{itemize}

\begin{tcolorbox}[breakable,colback=teal!4,colframe=teal!35!black,title={Before you move on}]
What does \ru{крокодил} mean? (\textquotedblleft{}A crocodile\textquotedblright{}.) Say it once more.
\end{tcolorbox}
\section[\texorpdfstring{\ru{жираф} (zhiráf)}{жираф (zhiráf)}]{\ru{жираф} (zhiráf) — a giraffe}

\label{lesson:RU-C225-zhiraf}



\begin{quote}
Before the new one: say the Russian for a sheep, then the Russian for a crocodile.
\end{quote}

\subsection*{You'll want to know: \ru{жираф}}
\textbf{\ru{жираф}} (\emph{zhiráf}) — \textquotedblleft{}a giraffe\textquotedblright{}.

\begin{cousinweb}
\textbf{\ru{жираф}} came through French \emph{girafe}, and the older \emph{\ru{жирафа}} through German or Italian, from Arabic \emph{zurāfa}, \textquotedblleft{}giraffe\textquotedblright{} (Vasmer).
\end{cousinweb}

\begin{grammarlens}[title={What you've built}]
One more word for this chapter.
\end{grammarlens}

\subsection*{Your turn}
\begin{itemize}
\raggedright
  \item \emph{Say it:} \emph{zhiráf}
  \item \emph{Say it:} \emph{zhiráf}, once more
  \item \emph{Say it:} say \emph{zhiráf}
  \item \emph{Recall:} say the Russian for a sheep, then the Russian for a crocodile, then say \emph{zhiráf} again
\end{itemize}

\begin{tcolorbox}[breakable,colback=teal!4,colframe=teal!35!black,title={Before you move on}]
What does \ru{жираф} mean? (\textquotedblleft{}A giraffe\textquotedblright{}.) Say it once more.
\end{tcolorbox}
\section[\texorpdfstring{\ru{зебра} (zébra)}{зебра (zébra)}]{\ru{зебра} (zébra) — a zebra}

\label{lesson:RU-C225-zebra}



\begin{quote}
Before the new one: say the Russian for a crocodile, then the Russian for a giraffe.
\end{quote}

\subsection*{You'll want to know: \ru{зебра}}
\textbf{\ru{зебра}} (\emph{zébra}) — \textquotedblleft{}a zebra\textquotedblright{}.

\begin{cousinweb}
\textbf{\ru{зебра}} more likely came through English or Dutch \emph{zebra} than straight from French \emph{zèbre}; Vasmer traces the word to an Angolan \emph{zerba}, and others to Latin.
\end{cousinweb}

\begin{grammarlens}[title={What you've built}]
One more word for this chapter.
\end{grammarlens}

\subsection*{Your turn}
\begin{itemize}
\raggedright
  \item \emph{Say it:} \emph{zébra}
  \item \emph{Say it:} \emph{zébra}, once more
  \item \emph{Say it:} say \emph{zébra}
  \item \emph{Recall:} say the Russian for a crocodile, then the Russian for a giraffe, then say \emph{zébra} again
\end{itemize}

\begin{tcolorbox}[breakable,colback=teal!4,colframe=teal!35!black,title={Before you move on}]
What does \ru{зебра} mean? (\textquotedblleft{}A zebra\textquotedblright{}.) Say it once more.
\end{tcolorbox}
\section[\texorpdfstring{\ru{верблюд} (verblyúd)}{верблюд (verblyúd)}]{\ru{верблюд} (verblyúd) — a camel}

\label{lesson:RU-C225-verblyud}



\begin{quote}
Before the new one: say the Russian for a giraffe, then the Russian for a zebra.
\end{quote}

\subsection*{You'll want to know: \ru{верблюд}}
\textbf{\ru{верблюд}} (\emph{verblyúd}) — \textquotedblleft{}a camel\textquotedblright{}.

\begin{cousinweb}
\textbf{\ru{верблюд}} — Old Russian \emph{velĭbludŭ}; an old loan from Gothic \emph{ulbandus}, reshaped after \emph{\ru{велий}}, \textquotedblleft{}great\textquotedblright{} (Vasmer). The Gothic word goes back to Greek \emph{eléphas}, \textquotedblleft{}elephant\textquotedblright{}: the name passed from one large beast to another.
\end{cousinweb}

\begin{grammarlens}[title={What you've built}]
One more word for this chapter.
\end{grammarlens}

\subsection*{Your turn}
\begin{itemize}
\raggedright
  \item \emph{Say it:} \emph{verblyúd}
  \item \emph{Say it:} \emph{verblyúd}, once more
  \item \emph{Say it:} say \emph{verblyúd}
  \item \emph{Recall:} say the Russian for a giraffe, then the Russian for a zebra, then say \emph{verblyúd} again
\end{itemize}

\begin{tcolorbox}[breakable,colback=teal!4,colframe=teal!35!black,title={Before you move on}]
What does \ru{верблюд} mean? (\textquotedblleft{}A camel\textquotedblright{}.) Say it once more.
\end{tcolorbox}
\section[\texorpdfstring{\ru{олень} (olén)}{олень (olén)}]{\ru{олень} (olén) — a deer}

\label{lesson:RU-C225-olen}



\begin{quote}
Before the new one: say the Russian for a zebra, then the Russian for a camel.
\end{quote}

\subsection*{You'll want to know: \ru{олень}}
\textbf{\ru{олень}} (\emph{olén}) — \textquotedblleft{}a deer\textquotedblright{}.

\begin{cousinweb}
\textbf{\ru{олень}} $\leftarrow$ Proto-Slavic \emph{*olenĭ} — inherited. Vasmer lists Lithuanian \emph{élnis}, \textquotedblleft{}deer, elk\textquotedblright{}, Greek \emph{élaphos}, \textquotedblleft{}deer\textquotedblright{}, and Welsh \emph{elain}, \textquotedblleft{}hind\textquotedblright{}.
\end{cousinweb}

\begin{grammarlens}[title={What you've built}]
One more word for this chapter.
\end{grammarlens}

\subsection*{Your turn}
\begin{itemize}
\raggedright
  \item \emph{Say it:} \emph{olén}
  \item \emph{Say it:} \emph{olén}, once more
  \item \emph{Say it:} say \emph{olén}
  \item \emph{Recall:} say the Russian for a zebra, then the Russian for a camel, then say \emph{olén} again
\end{itemize}

\begin{tcolorbox}[breakable,colback=teal!4,colframe=teal!35!black,title={Before you move on}]
What does \ru{олень} mean? (\textquotedblleft{}A deer\textquotedblright{}.) Say it once more.
\end{tcolorbox}
